
On
prend pour famille $(\xi_i)_{i\in I}$ la famille de tous les éléments de $K^{\bullet}(X)$. Les hypothèses de 1.2 sont alors vérifiées, et par suite l'homomorphisme :
\[
\varphi:\bigl(K_\bullet(S)\bigr)^{(K^{\bullet}(X))}\longrightarrow K_\bullet(X)
\]
est surjectif. D'après la définition de $\varphi$, ceci équivaut à dire que l'homomorphisme :
\[
K_\bullet(S)\otimes_{K^{\bullet}(S)}K^{\bullet}(X)\longrightarrow K_\bullet(X)
\]
est surjectif.

\textbf{Corollaire 1.4.} \uline{Soit} $f:X\to S$ \uline{un morphisme plat de schémas noethériens}. \uline{On suppose que pour tout point} $\underline{s}$ \uline{de} $S$, \uline{l'homomorphisme canonique} $\mathbb Z\to K_\bullet(X_s)$ \uline{est un isomorphisme}. \uline{Alors l'homomorphisme} :
\[
f^*:K_\bullet(S)\longrightarrow K_\bullet(X)
\]
\uline{est surjectif}. \uline{Si de plus} $f$ \uline{possède une section} $g$ \uline{de tor-dimension finie}, \uline{alors} $f^*$ \uline{est un isomorphisme}.

On prend pour famille $(\xi_i)_{i\in I}$ la famille réduite au seul élément $1\in K^{\bullet}(X)$; on a alors $\varphi=f^*$, et, les conditions de 1.2 étant remplies, $f^*$ est surjectif.

Si $f$ admet une section $g$ de tor-dimension finie, on peut définir un homomorphisme $g^*:K_\bullet(X)\to K_\bullet(S)$ (IV 2.12). Comme on a $g^*\circ f^*=(f\circ g)^*=\operatorname{Id}$ d'après IV 2.12.1.1, $f^*$ est injectif, donc est un isomorphisme.

\textbf{Remarque 1.5.} Pour que $g$ soit de tor-dimension finie, il faut et il suffit que les fibres de $X$ soient régulières aux points de $g(S)$, i.e.\ que $f$ soit régulier aux points de $g(S)$ (ou encore lisse en ces points si on suppose $f$
